% ============================================================================
%  附录　标准答案　—— 学生版讲义
%  只列有唯一标准答案的练习；开放性问题不列。
%  每一条答案都与 03-讲义/教师版/第NN课-*.tex 逐条核对过。
% ============================================================================

\clearpage
\phantomsection
\addcontentsline{toc}{chapter}{附录　标准答案}%
\renewcommand{\xhcurrentlesson}{附录　标准答案}%
\markboth{\xhcourse}{附录　标准答案}%
\vspace*{0.3cm}%
\noindent\colorbox{xhteal!12}{%
  \parbox[c][1.25cm][c]{\linewidth}{%
    \centering
    {\heiti\fontsize{14pt}{20pt}\selectfont\color{xhteal} 附录\quad\textbf{标准答案}}%
  }%
}\\[10pt]
\noindent{\small\songti\color{xhgray} 这里只列\textbf{有唯一答案}的题；「想一想」里让你自己举例、自己说想法的那类没有标准答案，所以我们不列。写完先自己对一遍，再来找老师。}

\section*{第 1 课　先导课：抽象代数是什么}
\begin{enumerate}
  \item \(999999\div7=142857\)——正好整除，一点不剩。
  \item \(142857\times7=999999\)。所以 \(999999\div7=142857\) 和 \(142857\times7=999999\) 是同一件事，正着说和倒着说。
\end{enumerate}

\section*{第 2 课　从余数到模运算：单位元与逆元}
\begin{enumerate}
  \item 今天是星期三，\(100\) 天以后是\textbf{星期五}（\(100=7\times14+2\)，扔掉 \(14\) 个整星期，只剩 \(2\) 天）。
  \item \(365\equiv1\pmod7\)，\(\quad1000\equiv6\pmod7\)。
  \item 手指回答（模 \(12\) 加法）：\(10+5\to3\)，\(\ 7+7\to2\)，\(\ 9+4\to1\)，\(\ 6+6\to0\)，\(\ 3+3+3+3+3\to3\)。
  \item 「取较大者」：\(3*5=5\)，\(\ (-2)*1=1\)。它是运算——任取两个有理数，结果还是有理数，而且答案唯一。
  \item 模 \(12\) 加法里的单位元是 \(0\)（走 \(0\) 格就是不动，\(a+0=a\)）。
  \item 单位元表：模 \(12\) 加法是 \(0\)（\(a+0=a\)）；有理数加法是 \(0\)（\(a+0=a\)）；有理数乘法是 \(1\)（\(a\times1=a\)）。一句话：加法的单位元是 \(0\)，乘法的单位元是 \(1\)。
  \item 逆元表：有理数加法——单位元 \(0\)，\(a\) 的逆元是相反数 \(-a\)；有理数乘法——单位元 \(1\)，\(a\) 的逆元是倒数 \(1/a\ (a\neq0)\)。
  \item 最后那张表：有理数加法 \(0\) 与 \(-a\)；有理数乘法 \(1\) 与 \(1/a\ (a\neq0)\)；模 \(12\) 加法 \(0\) 与 \(12-a\)（这一格第 2 课故意留空，第 3 课补上）。
\end{enumerate}

\section*{第 3 课　撤回一步：逆元}
\begin{enumerate}
  \item 模 \(12\) 加法里，\(3\to9\)，\(\ 7\to5\)，\(\ 5\to7\)。
  \item 搭档全表：\(0\to0,\ 1\to11,\ 2\to10,\ 3\to9,\ 4\to8,\ 5\to7,\ 6\to6,\ 7\to5,\ 8\to4,\ 9\to3,\ 10\to2,\ 11\to1\)。自己配自己的是 \(0\) 和 \(6\)。
  \item 乘 \(3\) 的转盘：\(1\to3,\ 2\to6,\ 3\to2,\ 4\to5,\ 5\to1,\ 6\to4\)。
  \item 第二行里没有两个数字是同一个（\(3,6,2,5,1,4\) 各出现一次）。所以「乘 \(3\)」把六个位置\textbf{重新排了一遍}。
  \item \(3x\bmod7\) 的表：\(3,\ 6,\ 2,\ 5,\ 1,\ 4\)。我找到 \(x=5\)（\(3\times5=15=2\times7+1\equiv1\pmod7\)）。
  \item 模 \(6\)：\(2x\bmod6\) 的表是 \(2,\ 4,\ 0,\ 2,\ 4,\ 0\)——找\textbf{不到} \(1\)。因为 \(2x\) 永远是偶数；更根本地说，\(2\) 和 \(6\) 有公因数 \(2\)。
  \item 对照表：有理数加法——\(0\) 与 \(-a\)；有理数乘法——\(1\) 与 \(1/a\ (a\neq0)\)；模 \(7\) 乘法——\(1\) 与表里的搭档（\(1\leftrightarrow1,\ 2\leftrightarrow4,\ 3\leftrightarrow5,\ 6\leftrightarrow6\)）。
  \item 模 \(7\) 乘法里，\(6\) 的逆元是 \(6\)——它和自己互逆（\(6\times6=36\equiv1\)）。
  \item 模 \(7\) 乘法里，\(2\) 的逆元是 \(4\)；\(4\) 的逆元是 \(2\)（\(2\times4=8\equiv1\)）。
\end{enumerate}

\section*{第 4 课　群的四条规则}
\begin{enumerate}
  \item 三个老熟人：模 \(12\) 加法——封闭 \(\checkmark\)、结合 \(\checkmark\)、单位元 \(0\)、逆元 \(12-a\)；有理数加法——\(\checkmark\)、\(\checkmark\)、\(0\)、\(-a\)；有理数乘法——\(\checkmark\)、\(\checkmark\)、\(1\)、\(1/a\ (a\neq0)\)。
  \item 有理数乘法里 \(0\) 没有逆元（要找 \(b\) 使 \(0\times b=1\)，而左边永远是 \(0\)）。少了一条逆元，四条凑不齐，所以有理数乘法\textbf{不是群}。
  \item \(3\oplus5=9\)，\(\ (-2)\oplus0=-1\)。
  \item 一条一条验：封闭 \(\checkmark\)（结果还是有理数）；结合 \(\checkmark\)（两边都等于 \(a+b+c+2\)）；单位元 \(e=-1\)（\textbf{不是} \(0\)）；逆元 \(b=-a-2\)。
  \item 手指回答：\(0\) 的逆元是 \(-2\)；\(1\) 的逆元是 \(-3\)；\(-3\) 的逆元是 \(1\)。
  \item \((1\oplus3)\oplus5=2\oplus5=3.5\)，\(\ 1\oplus(3\oplus5)=1\oplus4=2.5\)。两个结果\textbf{不一样}，所以它不满足结合律。
  \item 先甲后乙：\(1\to2\to3\)；先乙后甲：\(1\to2\to4\)。结果不一样（\(3\neq4\)），所以交换律不成立。
  \item 甲、乙、甲三个动作连着做：结果\textbf{一样}（先合前两个或先合后两个，都是「甲、乙、甲」）。
  \item 开头那题（看下面那张表的标题）：一共 \textbf{\(4\)} 条规则——封闭、结合、单位元、逆元；表的标题里\textbf{没有}「交换律」。
\end{enumerate}

\section*{第 5 课　三角形的对称与置换记号}
\begin{enumerate}
  \item 甲 \(=\) 乘 \(2\)，乙 \(=\) 加 \(1\)，取 \(x=1\)：先甲后乙得 \(3\)，先乙后甲得 \(4\)，不一样。
  \item 边看边记（只记 \(A\)、\(C\)）：不动——\(A\to A\)，\(C\to C\)；顺时针转 \(120^\circ\)——\(A\to B\)，\(C\to A\)；顺时针转 \(240^\circ\)——\(A\to C\)，\(C\to B\)。
  \item 一共数出 \(6\) 个动作：转动 \(3\) 个、翻折 \(3\) 个。
  \item 跟做：顺时针转 \(240^\circ\) 把 \(A,B,C\) 分别送到 \(C,A,B\)；沿过 \(B\) 的轴翻折把 \(A,B,C\) 分别送到 \(C,B,A\)。
  \item 六张两行式里，下排和上排一模一样的只有「不动」那一张（\(A,B,C\) 各自不动，也就是 \(A\to A\)、\(B\to B\)、\(C\to C\)）。
  \item 把 \(A,B,C\) 分别送到 \(B,A,C\) 的那个动作，是\textbf{沿过 \(C\) 的轴翻}（\(C\) 不动，\(A\) 与 \(B\) 对调）。做这个动作时，纸片\textbf{要翻面}。
\end{enumerate}

\section*{第 6 课　复合作为运算：\(S_3\) 是个群}
\begin{enumerate}
  \item 跟做：\(rf\)——\(A\to B\to C\)，\(B\to C\to B\)，\(C\to A\to A\)；\(fr\)——\(A\to A\to B\)，\(B\to C\to A\)，\(C\to B\to C\)。结论：\(rf\) 与 \(fr\) \textbf{不一样}（\(rf\neq fr\)）。
  \item 四条规则：封闭\textbf{成立}（\(rf\)、\(fr\) 仍是这 \(6\) 个之一）；结合\textbf{成立}；单位元是 \(e\)（不动）；逆元——\(r\) 的逆是 \(r^2\)，\(r^2\) 的逆是 \(r\)，\(f\) 的逆是 \(f\)。
  \item 乘法表：\(r\) 行是 \(r,\ r^2,\ e,\ rf,\ r^2f,\ f\)；\(f\) 行是 \(f,\ r^2f,\ rf,\ e,\ r^2,\ r\)；挑战行 \(r^2\) 行是 \(r^2,\ e,\ r,\ r^2f,\ f,\ rf\)。从 \(f\) 行的第二个格子还能读出 \(fr=r^2f\)。
  \item \(r\) 那一行里，没有动作出现两次，也没有动作没出现——每一行都是同样 \(6\) 个动作的\textbf{一次重排}。
\end{enumerate}

\section*{第 7 课　群里的小世界：阶与子群}
\begin{enumerate}
  \item 阶表：\(e\)（不动）连做 \(1,2,3\) 次都是 \(e\)，阶是 \(1\)；\(r\)（转 \(120^\circ\)）连做 \(1,2,3\) 次是 \(r,\ r^2,\ e\)，阶是 \(3\)；\(f\)（沿过 \(A\) 的轴翻）连做 \(1,2,3\) 次是 \(f,\ e,\ f\)，阶是 \(2\)。
  \item 挑战行 \(r^2\)（转 \(240^\circ\)）：连做 \(1,2,3\) 次是 \(r^2,\ r,\ e\)，所以阶是 \textbf{\(3\)}（不是 \(2\)）。
  \item 阶 \(2\) 的都是翻折，阶 \(3\) 的都是转动。
  \item \(\{e,r,r^2\}\)：填出的表里结果全都还在 \(e,r,r^2\) 里，没有跑出去，所以它\textbf{是}子群。
  \item \(\{e,r,f\}\)：\(r\cdot r=r^2\)、\(r\cdot f=fr=r^2f\)，这两个结果都不在 \(H\) 里（跑出去了），所以 \(H\) \textbf{不是}子群。
  \item \(S_3\) 里各子群的大小是 \(1,2,3,6\)，它们\textbf{都能整除} \(S_3\) 的大小 \(6\)（都是 \(6\) 的因数）。
\end{enumerate}

\section*{第 8 课　正方形的对称：\(D_4\) 的 \(8\) 个元素}
\begin{enumerate}
  \item 边看边记（只记 \(1\)、\(3\)）：不动——\(1\to1\)，\(3\to3\)；顺时针转 \(90^\circ\)——\(1\to2\)，\(3\to4\)；顺时针转 \(180^\circ\)——\(1\to3\)，\(3\to1\)；顺时针转 \(270^\circ\)——\(1\to4\)，\(3\to2\)。
  \item 把正方形转过 \(45^\circ\) 不算对称动作：转完图形与原来的位置\textbf{不重合}（顶点对不上顶点），和第 5 课三角形「歪着挪一下」是同一件事。
  \item 正方形一共有 \(8\) 个动作：转动 \(4\) 个、翻折 \(4\) 个。转 \(90^\circ\) 要连做 \(4\) 次才回到原位；三角形里转 \(120^\circ\) 连做 \(3\) 次。
  \item 乘法表：\(r\) 行是 \(r,\ r^2,\ r^3,\ e,\ rf,\ r^2f,\ r^3f,\ f\)；\(f\) 行是 \(f,\ r^3f,\ r^2f,\ rf,\ e,\ r^3,\ r^2,\ r\)；挑战行 \(r^2\) 行是 \(r^2,\ r^3,\ e,\ r,\ r^2f,\ r^3f,\ f,\ rf\)。从 \(f\) 行的第二个格子还能读出 \(fr=r^3f\)。
  \item \(r\) 那一行里没有动作出现两次——这一行是 \(8\) 个动作的\textbf{一次重排}。
\end{enumerate}

\section*{第 9 课　正方形里的小群：子群与二面体群}
\begin{enumerate}
  \item \(\{e,r^2\}\)：\(r^2\cdot r^2=e\)，没有跑出去，所以它\textbf{是}子群；\(\{e,r\}\)：\(r\cdot r=r^2\)，这个结果不在 \(H\) 里，所以它\textbf{不是}子群。
  \item \(H=\{e,r^2,f,r^2f\}\) 的表：所有格子的结果都落在这 \(4\) 个动作里，没有跑出去，所以它\textbf{是}子群，大小是 \(4\)。
  \item 子群大小 \(1,2,4,8\) 和 \(D_4\) 的大小 \(8\) 的关系：它们\textbf{都能整除} \(8\)（都是 \(8\) 的因数）。
  \item 左右两条路：两条路的终点\textbf{一样}（都是 \(A\to B\)、\(B\to A\)、\(C\to C\)），也就是 \(fr=r^3f\)。
  \item \(D_3\) 有 \(6\) 个动作；\(D_5\)（正五边形）有 \(10\) 个；\(D_{12}\) 有 \(24\) 个。规律：正 \(n\) 边形一共 \(2n\) 个动作。
\end{enumerate}

\section*{第 10 课　生成元、循环群与循环群的子群}
\begin{enumerate}
  \item 一直加 \(3\)：\(0\to3\to6\to9\to12\)（回到 \(0\)）。\(\langle3\rangle=\{0,3,6,9\}\)，一共 \(4\) 个。
  \item 一直加 \(5\)：\(0\to5\to10\to3\to8\to1\to6\to11\to4\to9\to2\to7\to0\)，一共走遍了 \(12\) 个数（\(5\) 是生成元，模 \(12\) 加法是循环群）。
  \item \(3\) 与 \(12\) 的公因数是 \(1,3\)（最大的是 \(3\)），\(12\div3=4\)；\(5\) 与 \(12\) 只有公因数 \(1\)，\(12\div1=12\)。公因数越大，圈住的越少。
  \item 模 \(12\) 各元素的阶：\(1\to12\)，\(2\to6\)，\(3\to4\)，\(4\to3\)，\(5\to12\)，\(6\to2\)，\(7\to12\)，\(8\to3\)，\(9\to4\)，\(10\to6\)，\(11\to12\)，\(0\to1\)。
  \item \(\langle4\rangle\) 和 \(\langle8\rangle\) \textbf{是}同一个子群，都是 \(\{0,4,8\}\)。
  \item \(1,2,3,4,6,12\) 和 \(12\) 的关系：它们\textbf{都能整除} \(12\)。\(S_3\) 里是 \(1,2,3,6\)，\(D_4\) 里是 \(1,2,4,8\)——三次的结论一样：子群的大小都整除群的大小。
\end{enumerate}

\section*{第 11 课　幂的尾巴转圈了：从现象到猜想}
\begin{enumerate}
  \item \(2^n\) 除以 \(7\) 的余数在 \(2,4,1\) 里打转。第 \(7\) 个余数是 \(2\)，和第 \(1\) 个一样；第 \(2026\) 个也是 \(2\)（\(2026=3\times675+1\)，落在循环里的第 \(1\) 个）。
  \item \(2^n\) 除以 \(5\)：\(2,\ 4,\ 3,\ 1\)——第 \(4\) 步第一次到 \(1\)。
  \item \(3^n\) 除以 \(5\)：\(3,\ 4,\ 2,\ 1\)——第 \(4\) 步第一次到 \(1\)。
  \item \(3^n\) 除以 \(7\)：\(3,\ 2,\ 6,\ 4,\ 5,\ 1\)——第 \(6\) 步第一次到 \(1\)。
  \item 归纳两句话：每一列的步数\textbf{都整除} \(p-1\)（\(p=5\) 时整除 \(4\)，\(p=7\) 时整除 \(6\)）；走满 \(p-1\) 步时余数都是 \(1\)（\(2^4\)、\(3^4\) 除以 \(5\) 各余 \(1\)，\(2^6\)、\(3^6\) 除以 \(7\) 各余 \(1\)）。
  \item 「试出来的」\textbf{不算数}——例子再多也只是例子，必须证明（下一课就证）。
\end{enumerate}

\section*{第 12 课　费马小定理的证明}
\begin{enumerate}
  \item 取余数那一行：\(3,\ 6,\ 2,\ 5,\ 1,\ 4\)。
  \item 它是 \(1,2,3,4,5,6\) 的一次\textbf{重排}；没有哪个数字出现两次，也没有哪个数字没出现。
  \item \(3\times6\times9\times12\times15\times18=3^6\times(1\times2\times3\times4\times5\times6)=3^6\times6!\)。
  \item \(6!\) 里没有一个因数是 \(7\)，所以它\textbf{有逆元}，两边可以同时划掉。划完得 \(3^6\equiv1\pmod7\)（\(3^6=729=104\times7+1\)，余 \(1\)）。
  \item 换成字母：余数是 \(1\text{—}(p-1)\) 的一次重排；\(a^{p-1}\times(p-1)!\equiv(p-1)!\)；最后 \(a^{p-1}\equiv1\pmod p\)。
  \item 两块砖：砖一——\(a,2a,\dots,(p-1)a\) 的余数正好是 \(1,2,\dots,p-1\) 的一次重排；砖二——模 \(p\) 里只要不是 \(p\) 的倍数就有逆元，所以两边的 \((p-1)!\) 可以划掉。
  \item 证明的四步：①乘 \(a\) 再取余，是 \(1,\dots,p-1\) 的一次重排；②两边各自相乘；③划掉都有的 \((p-1)!\)（它有逆元）；④得到 \(a^{p-1}\equiv1\pmod p\)。
\end{enumerate}

\section*{第 13 课　它有什么用}
\begin{enumerate}
  \item \(999999=10^6-1\)；由费马小定理 \(10^6\equiv1\pmod7\)；代回去 \(10^6-1\equiv1-1=0\pmod7\)，所以 \(7\) 整除 \(999999\)。
  \item 余数是 \(0\)，说明 \(999999\) 能被 \(7\) 整除。从头到尾我们\textbf{没有}做过一次除法，也\textbf{没有}用过计算器。
  \item \(200=10\times20\)，\(\ 3^{200}=(3^{10})^{20}\equiv1^{20}=1\pmod{11}\)。
  \item \(5^{2026}\) 除以 \(7\)：第一步先算 \(2026\div6\) 的余数，得 \(4\)；答案与 \(5^4\equiv2\) 相同，所以余数是 \(2\)。
  \item 身份证示例（前 \(17\) 位 \(35060420130520123\)）：加权和是 \(218\)；\(218=11\times19+9\)；余数是 \(9\)，查表得校验码是 \(3\)。完整号码 \(350604201305201233\)。
  \item 改掉前 \(17\) 位里的任意一位，加权和就变了，余数跟着变，校验码对不上——所以机器能发现号码不对。
  \item 用你自己的号码验一遍（分三段填小计）：第 \(1\)—\(6\) 位、第 \(7\)—\(14\) 位、第 \(15\)—\(17\) 位分别求小计，再加起来。以示例号 \(350604201305201233\) 为例，三段小计是 \textbf{\(112\)}、\textbf{\(84\)}、\textbf{\(22\)}，合计 \(218\)。
    （前六位和你一样的同学，第一段直接用 \(112\)。）
  \item 选做（只算 \(9\) 项）：前 \(6\) 位 \(+\) 后 \(3\) 位 \(=112+22=134\)，再用已知总和倒推中间 \(8\) 项：\(218-134=\)\textbf{\(84\)}。
\end{enumerate}

\section*{第 14 课　复习、总结与展望}
\begin{enumerate}
  \item 三条线为什么会汇到一起：它们共同满足群的\textbf{四条规则}——封闭、结合、单位元、逆元。
  \item 合上书抽查：①模 \(12\) 加法里 \(5\) 的逆元是 \(7\)；②正三角形有 \(6\) 个对称动作，正方形有 \(8\) 个；③\(rf\) 和 \(fr\) \textbf{不一样}；④\(2^{100}\) 除以 \(7\) 余 \(2\)。
  \item 四题各属哪一支：①左（模 \(12\) 加法）；②右（图形对称）；③右（图形对称）；④下（幂的余数）。
  \item 为什么说「整数是环，不是域」：整数里不是每个非零元素都能做除法——例如 \(2\) 的倒数 \(1/2\) 不在整数里。
\end{enumerate}
